\section{Lecture 3 (January 15, 2025): Affine groups and representations}\label{lec3}
\noindent\textit{Scribe: Andrew Aguilar}

\subsection{Recovering the Hopf structure}
We finish the antiequivalence between affine group schemes and commutative
Hopf algebras begun in Lecture~\ref{lec2}.

\begin{theorem}
The coordinate-algebra construction gives an antiequivalence between
$\mathsf{AffGrp}_k$ and $\mathsf{CommHopf}_k$.
\end{theorem}
\begin{proof}
The two constructions of the previous lecture agree with the underlying
Yoneda antiequivalence. It remains to check the additional structures.
Let $A$ be a commutative Hopf algebra and let
$G(R)=\Hom_{k\text{-}\mathsf{Alg}}(A,R)$ with its convolution product.
The two algebra maps
\[
 \iota_l(a)=a\otimes1,\qquad \iota_r(a)=1\otimes a
\]
are points of $G(A\otimes A)$. Their product is
\[
 (\iota_l\iota_r)(a)
   =\sum(a_1\otimes1)(1\otimes a_2)
   =\sum a_1\otimes a_2=\Delta(a).
\]
Thus the group law recovers the original comultiplication. The identity
of $G(k)$ is $\epsilon:A\to k$, and the inverse of the universal point
$\id_A\in G(A)$ is $S:A\to A$. This recovers the counit and antipode.
Conversely, Yoneda identifies the natural group operations on any affine
$G$ with precisely these three algebra maps, so they are recovered as well.
Finally, a natural map of affine schemes is a group homomorphism exactly
when its coordinate-algebra map commutes with these operations. Hence the
constructions are inverse on objects and morphisms, up to their natural
Yoneda identifications.
\end{proof}

For clarity, the multiplication on $A\otimes B$ used above is
\[
 (a\otimes b)(a'\otimes b')=aa'\otimes bb'.
\]
It is the composite of $\id_A\otimes\tau_{B,A}\otimes\id_B$ with
$m_A\otimes m_B$, where $\tau_{B,A}(b\otimes a')=a'\otimes b$.

\begin{example}
The additive group $\mathbb G_a$ has coordinate algebra $k[x]$, with
\[
 \Delta(x)=x\otimes1+1\otimes x,\qquad
 \epsilon(x)=0,\qquad S(x)=-x.
\]
These formulas express addition, zero and negation.
\end{example}

\begin{example}
For $\GL_n$, put $X=(x_{ij})$ and $d=\det X$. Its coordinate algebra is
$A=k[x_{ij},d^{-1}]$. Matrix multiplication gives
\[
 \Delta(x_{ij})=\sum_{l=1}^n x_{il}\otimes x_{lj},\qquad
 \Delta(d^{-1})=d^{-1}\otimes d^{-1}.
\]
The identity matrix and inverse matrix give
\[
 \epsilon(x_{ij})=\delta_{ij},\qquad S(x_{ij})=(X^{-1})_{ij},
 \qquad \epsilon(d^{-1})=1,\quad S(d^{-1})=d.
\]
\end{example}
\begin{exercise}
Check directly that these formulas respect $dd^{-1}=1$ and satisfy the
counit and antipode identities.
\end{exercise}

\begin{example}
For $\mathbb G_m=\GL_1$ the algebra is $k[x,x^{-1}]$, and
\[
 \Delta(x)=x\otimes x,\qquad\epsilon(x)=1,\qquad S(x)=x^{-1}.
\]
\end{example}

\subsection{Rational representations and comodules}
Fix a field $k$, an affine group scheme $G$ and a $k$-vector space $V$.
Write $V_R=V\otimes_kR$ and
$\GL(V)(R)=\Aut_R(V_R)$. When $V$ is finite-dimensional, the functor
$R\mapsto V_R$ is an affine scheme represented by $\Sym(V^*)$, and
$\GL(V)$ is an affine group scheme isomorphic to $\GL_{\dim V}$ after
choosing a basis. For arbitrary $V$, we use these functors without asserting
that they are affine schemes.

\begin{definition}\label{RatRepDef}
A \emph{rational $G$-module} is a vector space $V$ with an action
$G(R)\curvearrowright V_R$ by $R$-linear automorphisms, natural in every
commutative $k$-algebra $R$. Equivalently, it is a morphism of group-valued
functors $G\to\GL(V)$.
\end{definition}

\begin{definition}\label{ComodDef}
Let $(A,\Delta,\epsilon)$ be a coalgebra over $k$. A \emph{right
$A$-comodule} is a vector space $M$ with a linear map
$\rho:M\to M\otimes A$ such that
\[
 (\rho\otimes\id_A)\rho=(\id_M\otimes\Delta)\rho,
 \qquad(\id_M\otimes\epsilon)\rho=\id_M.
\]
These are coassociativity and counitality. In Sweedler notation we write
$\rho(m)=\sum m_0\otimes m_1$, with the summation understood to be finite
for each $m$.
\end{definition}

\begin{theorem}\label{rep-comodule-dictionary}
Let $A=k[G]$. Giving a rational $G$-module structure on $V$ is equivalent
to giving a right $A$-comodule structure on $V$. If $V$ is finite-dimensional,
these are also equivalent to a morphism of affine group schemes
$G\to\GL(V)$.
\end{theorem}
We prove this statement and its compatibility with morphisms in the next
lecture.
